\section{Research Objectives and Questions}
\label{sec:introduction_research_objectives}

The objective of this thesis is to investigate how different representation strategies in VR indoor rowing influence embodiment-related and experiential outcomes. Specifically, the thesis compares an ergometer-congruent representation with a boat-centric representation. The comparison is guided by the assumption that both strategies can be justified by different design goals: the ergometer-congruent mode may support sensorimotor alignment and agency, while the boat-centric mode may support perceived rowing realism and the impression of rowing on water.

The first objective is to design and implement a VR rowing prototype that can represent the same physical ergometer activity in two different ways. In the ergometer-congruent mode, the system emphasizes direct correspondence between the user's real movement and the virtual body or handle. In the boat-centric mode, the system emphasizes a rowing scene with boat and oar depiction, even if this requires transforming the user's handle movement into a more sport-like visual movement.

The second objective is to use this prototype as a research instrument. The system is not only an application prototype, but also a means of examining how physical-to-virtual movement mapping affects user experience. This requires the two modes to be comparable in as many respects as possible, so that the central difference remains the relation between the real ergometer movement and the virtual rowing depiction.

The third objective is to evaluate the two representation modes with respect to embodiment-related constructs. These include sense of agency, body ownership, self-location, and overall sense of embodiment. In addition, the study considers perceived rowing realism, comfort, and preference as exploratory design-facing outcomes.

The thesis is guided by the following main research question:

\begin{quote}
How does the trade-off between visuomotor congruence and realistic sport depiction affect embodiment, perceived rowing realism, and user preference in VR indoor rowing?
\end{quote}

This main question is divided into the following subquestions:

\begin{enumerate}
    \item How can an ergometer-congruent and a boat-centric rowing representation be implemented while keeping the physical task and virtual rowing context broadly comparable?

    \item How do the two representation modes differ in embodiment-related experience, particularly sense of agency, body ownership, self-location, and overall sense of embodiment?

    \item How do the two representation modes differ in perceived rowing realism, comfort, and user preference?

    \item What do the results suggest for the design of XR sport and exercise systems in which physical equipment and virtual sport actions do not fully coincide?
\end{enumerate}

These questions are intentionally framed as comparative and exploratory rather than as claims of expected superiority. It is plausible that the ergometer-congruent mode supports stronger agency because the visible movement more closely matches the user's physical action. It is also plausible that the boat-centric mode supports a stronger sense of rowing realism because it better represents the visual and cultural form of the sport. The study therefore examines how these dimensions interact rather than assuming that one form of fidelity is always preferable.
